\section{Conclusiones}\label{sec:conclusiones}

Construimos un modelo sencillo para estimar el rendimiento de 59 parcelas de cebada de temporal. Usa dos datos, el
régimen de registro del municipio y el rendimiento de las parcelas vecinas, combinados con un proceso gaussiano.
Esperamos un RMSE de 0.46 a 0.49\,t/ha en las 59 parcelas.

Cada pieza entró porque mejoró la validación de forma consistente, y cada idea que quedó fuera lo hizo porque no la
mejoró. Así quedaron fuera las variables satelitales y climáticas, que no aportaron una vez conocidos el régimen y los
vecinos.

Nuestros resultados sugieren que el límite de este reto está más en las etiquetas que en el algoritmo. El modelo mejora
poco a la media por régimen, y el redondeo y la fuente mixta del rendimiento pueden agregar un ruido que no pudimos
medir. Conocer la fuente de cada registro, la fecha de siembra o tener cortes de cosecha medidos permitiría comprobarlo,
y podría ayudar más que un modelo más complejo.

Sobre las preguntas guía, el NDVI y la lluvia pesan solo cuando falta el régimen, porque indican la zona. Con el régimen
y los vecinos, el NDDI es el índice que más aporta, pero su efecto es muy pequeño. La mejor combinación es régimen más
rezago espacial.

Para un productor, el error equivale a unos 2,500 pesos por hectárea. La estimación sirve para ordenar parcelas y
orientar la revisión de un crédito, junto con el juicio del analista.

Como trabajo futuro proponemos validar con el rendimiento real de las 59 parcelas cuando FIRA lo publique, sumar más
temporadas para estimar antes de la cosecha y probar la capa de crédito con datos reales.
